% ch13.tex — Lecture 13: The Bakery Algorithm (13 February 2026)
% Scribe: Group 17

% Local theorem/lemma/proof environments (amsthm is not loaded in main.tex).
\newtheorem{theorem}{Theorem}[chapter]
\newtheorem{lemma}[theorem]{Lemma}
\newenvironment{proof}[1][Proof]{%
  \par\noindent\textit{#1.}\enspace\ignorespaces
}{%
  \hfill$\square$\par\medskip
}

\chapter{Lecture 13: The Bakery Algorithm}

\section{Recap: Where We Left Off}

In the previous lecture, we studied Peterson's algorithm for 2-thread mutual
exclusion. We showed that it satisfies mutual exclusion, deadlock-freedom,
starvation-freedom, and bounded waiting (with $r = 1$). We then raised the
question: can we solve mutual exclusion for $n$ threads? The answer is the
\ppterm{Bakery Algorithm}, which we introduced briefly at the end of Lecture~12.
This lecture picks up from there and focuses on proving its correctness.

\section{The Bakery Algorithm}

The Bakery Algorithm, due to Lamport (1974), is a classical $n$-thread mutual
exclusion protocol. The intuition is simple: it works like a bakery where
customers take numbered tickets and are served in order. Each thread picks a
label (ticket number) and waits until it holds the smallest label among all
interested threads.

\subsection{Data Structures and Code}

The algorithm maintains two shared arrays:
\begin{itemize}[nosep]
  \item \ppinline{flag[0..n-1]}: a boolean array indicating whether thread $i$
        is interested in entering the critical section.
  \item \ppinline{label[0..n-1]}: an integer array storing the ticket number
        (label) of each thread.
\end{itemize}

The pseudocode is as follows:
\clearpage
\begin{ppprogram}{Bakery Algorithm (Herlihy-Shavit formulation)}{prog:bakery}
\begin{lstlisting}[language=Java]
class Bakery implements Lock {
  boolean[] flag;
  Label[]   label;

  public Bakery(int n) {
    flag  = new boolean[n];
    label = new Label[n];
    for (int i = 0; i < n; i++) {
      flag[i] = false; label[i] = 0;
    }
  }

  public void lock() {
    flag[i]  = true;                               // (D1)
    label[i] = max(label[0], ..., label[n-1]) + 1; // (D2)
    while (for some k:                              // (W)
      flag[k] && (label[i], i) > (label[k], k));
  }

  public void unlock() {
    flag[i] = false;
  }
}
\end{lstlisting}
\end{ppprogram}

\subsection{Structure of \ppinline{lock()}}

The \ppinline{lock()} method naturally splits into two phases:

\begin{enumerate}[nosep]
  \item \textbf{Doorway} (finite steps, lines D1--D2): The thread sets its flag
        and computes its label. This always terminates in a bounded number of
        steps.
  \item \textbf{Waiting} (potentially unbounded, line W): The thread spins in
        the \ppinline{while} loop until no other interested thread holds a
        lexicographically smaller (\ppinline{label}, id) pair.
\end{enumerate}

The waiting condition checks, for every other thread $k$: is $k$ interested
(\ppinline{flag[k]} is true) \emph{and} does $k$ have a smaller label (breaking
ties by thread id)? The comparison $(l_1,\mathit{id}_1) < (l_2,\mathit{id}_2)$
is lexicographic: first compare labels, and if equal, compare thread ids.

More precisely, thread $i$ exits the while loop and enters the CS only when,
for \emph{every} thread $k$:
\[
  \neg\,\texttt{flag}[k]
  \quad \lor \quad
  (\texttt{label}[i],\, i) \;\leq\; (\texttt{label}[k],\, k).
\]
That is, every other thread is either not interested, or has a label at least
as large as $i$'s (with tie-breaking by id).

\begin{ppnote}
\par
The version presented in class follows Herlihy and Shavit's textbook
formulation. Lamport's original 1974 algorithm uses a separate
\ppinline{choosing[i]} boolean variable (set to \ppinline{true} while thread
$i$ is computing its label, then reset to \ppinline{false}), and other threads
explicitly wait with \ppinline{while(choosing[j])} before inspecting $j$'s
label. Lamport's original also resets \ppinline{number[i] = 0} on unlock,
rather than using a separate flag.

In the Herlihy-Shavit version, \ppinline{flag[i]} is set to \ppinline{true}
\emph{before} the label is computed. This means any thread that sees
\ppinline{flag[i] = true} knows that thread $i$ is either still computing its
label or has already finished---in either case it is ``in the game,'' and other
threads account for it during the while-loop scan. The correctness argument
below relies on this write ordering and on sequential consistency (SC).
\end{ppnote}

\section{Does the Bakery Algorithm Guarantee Mutual Exclusion?}

This was the central question of today's lecture. We worked through it carefully.

\subsection{A Concrete Trace}

To build intuition, consider two threads $T_0$ and $T_1$ whose doorways
overlap so that both end up choosing the same label value. (If $T_0$'s
entire doorway completed before $T_1$ began, Lemma~\ref{lem:label-order}
would force $T_1$ to choose label 2, not 1. The interesting case is when
both threads' max-scans run before either label has been written.)

\begin{enumerate}[nosep]
  \item $T_0$ sets \ppinline{flag[0] = true} (line~D1) and starts its
        max-scan; it reads \ppinline{label[1] = 0} (the initial value).
  \item $T_1$ sets \ppinline{flag[1] = true} and starts its max-scan; it
        reads \ppinline{label[0] = 0} (still the initial value, since $T_0$
        has not reached line~D2 yet).
  \item $T_0$ finishes its scan with $\max=0$ and writes
        \ppinline{label[0] = 1}.
  \item $T_1$ finishes its scan with $\max=0$ and writes
        \ppinline{label[1] = 1}.
  \item $T_0$ enters the while loop and checks $k=1$: it reads
        \ppinline{flag[1] = true} and evaluates $(1, 0) > (1, 1)$, which is
        \textbf{false} (since $0 < 1$). The blocking conjunction is false,
        so $T_0$ enters the CS.
  \item $T_1$ enters the while loop and checks $k=0$: it reads
        \ppinline{flag[0] = true} and evaluates $(1, 1) > (1, 0)$, which is
        \textbf{true}. The blocking conjunction is true, so $T_1$ waits.
  \item $T_0$ finishes its critical section and sets
        \ppinline{flag[0] = false}.
  \item $T_1$ re-checks: \ppinline{flag[0]} is now \texttt{false}. The
        conjunction is false, so $T_1$ exits the loop and enters the CS.
\end{enumerate}

Even when both threads end up with the same label value, tie-breaking by
thread id ensures a total order and prevents simultaneous CS entry.

\subsection{Handling Concurrent Label Assignment}

\begin{ppkey}
\par
The label assignment is \emph{not} atomic. Two threads can concurrently compute
the same label value. This is fine: the lexicographic comparison $(\text{label},
\text{id})$ always breaks ties deterministically, since thread ids are unique.
\end{ppkey}

For instance, $P_1$ and $P_2$ might both compute label 1. Since $(1,1) < (1,2)$,
there is always a strict total order among (label, id) pairs.

More generally, suppose $P_1$ enters the CS first. Then from $P_1$'s perspective,
it observed $(k_1, 1) < (k_2, 2)$ for some label values $k_1, k_2$. From
$P_2$'s perspective, it needs $(m_2, 2) < (m_1, 1)$ to enter. Note that
$k_1, k_2$ and $m_1, m_2$ need not be the same values: these are the label
values as seen at different points in time. The key insight is that the
lexicographic ordering, combined with the structure of the algorithm, still
prevents simultaneous access.

\section{Formal Proof of Mutual Exclusion}

We now establish the correctness of the Bakery Algorithm rigorously. The proof
below follows the lecture, with some steps made more explicit after
cross-checking with standard references (Herlihy \& Shavit's textbook,
Lamport's original paper, Garg's textbook, and Shankar's notes).

Throughout, we assume \textbf{sequential consistency} (SC): all memory
operations appear to execute in some global total order consistent with the
program order of each thread, and every read returns the value of the most
recent write in that total order.

\subsection{Lemma: Doorway Order Implies Label Order}

To turn doorway timing into label arithmetic, we need a hypothesis strong
enough that $P_j$'s entire max-scan---not just its label-write---happens after
$P_i$'s label is in place. The right precondition is the ordering of the
\emph{doorway} of one thread relative to the \emph{start of the doorway} of
the other (this is the form Lamport states in his original 1974 paper, and is
the FCFS lemma in Herlihy and Shavit).

\begin{lemma}\label{lem:label-order}
If thread $P_i$'s doorway (lines D1 and D2) completes strictly before thread
$P_j$'s doorway begins, then $L_j > L_i$, and hence
$(L_j, j) > (L_i, i)$.
\end{lemma}

\begin{proof}
By hypothesis, $P_j$'s doorway begins after $P_i$'s doorway has completed. So
every step of $P_j$'s doorway, including every read in its max-scan, happens
after $P_i$'s write of $\texttt{label}[i] = L_i$ in line~(D2). Under SC, $P_j$'s
read of $\texttt{label}[i]$ therefore returns $L_i$ (the most recent write,
since $P_i$ writes $\texttt{label}[i]$ exactly once per lock acquisition and
has not yet started a new acquisition). Hence
\[
  L_j \;=\; \max(\ldots,L_i,\ldots) + 1 \;\geq\; L_i + 1 \;>\; L_i,
\]
which gives $(L_j, j) > (L_i, i)$ regardless of the ids.
\end{proof}

\begin{ppnote}
\par
The weaker hypothesis ``\emph{$P_i$ completes its label-write before $P_j$
completes its label-write}'' is \emph{not} enough on its own: $P_j$'s
max-scan can race in the middle of $P_i$'s doorway and read a stale
$\texttt{label}[i]$, producing $L_j \leq L_i$ even though $P_i$'s write
finishes first. This is why the lemma talks about doorways, not label
writes. The mutual exclusion proof below is structured to use only the
doorway form.
\end{ppnote}

\subsection{Theorem: Mutual Exclusion is Guaranteed}

\begin{theorem}\label{thm:me}
The Bakery Algorithm satisfies mutual exclusion.
\end{theorem}

\begin{proof}
Suppose for contradiction that, at some time $T$, threads $P_i$ and $P_j$ are
simultaneously in the CS. We consider one specific lock acquisition for each
of them (the one that placed them in the CS at time $T$) and write $L_i$,
$L_j$ for the labels chosen by those acquisitions.

\medskip
\noindent\textbf{Step 1: Establish a label ordering (WLOG).}
Since thread ids are distinct, $(L_i, i) \neq (L_j, j)$. Without loss of
generality, assume $(L_i, i) < (L_j, j)$. The plan is to derive a
contradiction by examining what $P_j$ must have observed in its waiting loop.

\medskip
\noindent\textbf{Step 2: $P_j$ must have observed $\texttt{flag}[i] = \texttt{false}$.}
For $P_j$ to have exited the while loop and entered the CS, it must have
found, at some final test $k=i$, that the blocking condition is false:
\[
  \neg\,\texttt{flag}[i] \;\;\lor\;\; (L_j, j) \,\leq\, (\texttt{label}[i], i).
\]
The second disjunct, were $P_j$ to read the value $L_i$, would give
$(L_j, j) \leq (L_i, i)$, contradicting Step~1. The same disjunct is
false for any value $\ell \leq L_i$ that $P_j$ could read in
$\texttt{label}[i]$, since then $(L_j, j) \leq (\ell, i) \leq (L_i, i)$
again contradicts Step~1. (And $\ell \leq L_i$ for every value $P_j$ can
read: in the Herlihy-Shavit formulation, $\texttt{label}[i]$ is never
reset, so the writes by $P_i$'s successive acquisitions form a strictly
increasing sequence; the current acquisition's $L_i$ is the maximum
written so far, and any earlier read returns either $0$ or a strictly
smaller previous label.) So the disjunct that must hold is the first
one: $P_j$ observed $\texttt{flag}[i] = \texttt{false}$ at its final
test for $k=i$.

\medskip
\noindent\textbf{Step 3: $P_j$'s observation forces $D_j \to D_i$.}
Because $P_i$ is in the CS at time $T$, it has not executed its
\ppinline{unlock()} yet. So between the moment $P_i$ wrote
$\texttt{flag}[i] = \texttt{true}$ in line~(D1) and time $T$, the value of
$\texttt{flag}[i]$ is continuously \ppinline{true}. The only way for $P_j$ to
have read $\texttt{flag}[i] = \texttt{false}$ is for $P_j$'s read to have
happened \emph{before} $P_i$'s write of $\texttt{flag}[i] = \texttt{true}$,
i.e., before $P_i$'s doorway began.

But $P_j$'s read of $\texttt{flag}[i]$ sits in $P_j$'s waiting loop, which
executes only after $P_j$'s doorway has completed. Combining,
\[
  (\text{$P_j$'s doorway completes}) \;\to\; (\text{$P_j$'s read})
  \;\to\; (\text{$P_i$'s doorway begins}),
\]
i.e., $D_j \to D_i$: $P_j$'s doorway ends strictly before $P_i$'s begins.

\medskip
\noindent\textbf{Step 4: Apply the Lemma to obtain a contradiction.}
By Lemma~\ref{lem:label-order} applied with $P_j$ as the earlier thread and
$P_i$ as the later thread, $L_i > L_j$, hence $(L_i, i) > (L_j, j)$. But
Step~1 assumed $(L_i, i) < (L_j, j)$, contradiction.
\end{proof}

\begin{ppnote}
\par
The slides express the contradiction more compactly as
``$(L_j, j) > (L_i, i) \wedge (L_j, j) < (L_i, i)$.'' The first conjunct is
the WLOG assumption from Step~1; the second comes from $P_j$'s observation,
once we have argued (Step~3) that $P_j$ cannot in fact have seen
$\texttt{flag}[i] = \texttt{false}$ \emph{after} $P_i$'s flag-write. The
critical content of the proof is therefore Step~3, which the slides leave
implicit.
\end{ppnote}

\subsection{Why the Proof Works: Three Key Ingredients}

It is worth stepping back and noting what makes the argument go through:

\begin{enumerate}[nosep]
  \item \textbf{Lexicographic ordering gives a total order.} Unique thread ids
        ensure no two threads ever share the same (label, id) pair, even when
        their label values are equal. This guarantees a well-defined ``winner''
        in every pairwise comparison and is what licenses the WLOG step.
  \item \textbf{Disjoint doorways force label growth}
        (Lemma~\ref{lem:label-order}). Under SC, if one doorway finishes
        before another begins, the later thread's max-scan must read the
        earlier thread's freshly-written label and choose a strictly larger
        label. This is the bridge from real-time ordering of doorways to the
        algorithm's logical ordering of labels.
  \item \textbf{The flag persists from doorway through CS.}
        \ppinline{flag[i]} is set in line~(D1) and only cleared on
        \ppinline{unlock()}. So while $P_i$ is in the CS, its flag is
        continuously \ppinline{true}; any other thread that observes
        $\texttt{flag}[i] = \texttt{false}$ must have read it \emph{before}
        $P_i$'s doorway began. This is what lets us turn $P_j$'s
        ``\texttt{flag}$[i] = \texttt{false}$'' observation into the
        doorway ordering $D_j \to D_i$.
\end{enumerate}

If any of these three ingredients were missing, the proof would break.

\subsection{Why \ppinline{flag[i]} is Set Before the Label (and Why It Matters)}

A natural question is: what if the order of lines (D1) and (D2) were swapped,
so we compute the label first and then set the flag? This breaks mutual
exclusion outright.

\begin{pppitfall}
\par
With the doorway ``label first, flag second,'' another thread can scan, see
$\texttt{flag}[i] = \texttt{false}$ during the gap between $i$'s label-write
and $i$'s flag-write, and enter the CS even though $i$ has already chosen its
label. Step~3 of the proof above silently relies on the flag being set before
the label.
\end{pppitfall}

\noindent
A concrete two-thread violation. Initially $\texttt{flag}[\cdot] = \texttt{false}$
and $\texttt{label}[\cdot] = 0$. Both threads use the swapped doorway, and ids
satisfy $i < j$.

\begin{enumerate}[nosep,leftmargin=2.2em]
  \item $P_i$ reads $\max(\ldots) = 0$ in its (now leading) max-scan and
        writes $\texttt{label}[i] = 1$. (\textit{$P_i$'s flag is still
        \texttt{false}}.)
  \item $P_j$ reads $\max(\ldots) = 1$ (it sees $i$'s new label) and writes
        $\texttt{label}[j] = 2$.
  \item $P_j$ writes $\texttt{flag}[j] = \texttt{true}$ and enters the while
        loop. For $k=i$, it reads $\texttt{flag}[i] = \texttt{false}$
        (\textit{$P_i$ has not reached its flag-write yet}); the conjunction is
        false, so $P_j$ skips $i$. $P_j$ has no other threads to wait on, so
        $P_j$ enters the CS.
  \item $P_i$ writes $\texttt{flag}[i] = \texttt{true}$ and enters the while
        loop. For $k=j$ it reads $\texttt{flag}[j] = \texttt{true}$ and
        $\texttt{label}[j] = 2$, and checks $(L_i, i) > (L_j, j)$, i.e.,
        $(1, i) > (2, j)$, which is false. $P_i$ also enters the CS.
\end{enumerate}

Both $P_i$ and $P_j$ are now in the CS at the same time. The bug is that the
gap between (the swapped) label-write and flag-write is observable: $P_j$ saw
$i$ as ``not interested'' even though $i$ was already mid-doorway. With the
correct order (D1) before (D2), $P_i$'s flag is set before $i$'s label is
written, so any other thread that sees $i$'s label or scans during $i$'s
doorway necessarily also sees $\texttt{flag}[i] = \texttt{true}$, and Step~3
of the mutual-exclusion proof goes through.

\section{Deadlock-Freedom}

\begin{ppkey}
\par
At any point, the interested thread with the smallest (label, id) pair cannot
be blocked by anyone: for every other thread $k$, either \ppinline{flag[k]} is
false or $(\texttt{label}[k], k) \geq (\texttt{label}[i], i)$. So it always
enters the CS, and deadlock is impossible.
\end{ppkey}

Consider any moment at which some thread is interested in entering the CS.
Each interested thread is in one of two states: still inside its doorway
(executing the bounded sequence (D1)--(D2)) or in the spin loop. Doorways are
straight-line code with no waiting, so any thread in its doorway will, in a
finite number of its own steps, transition into the spin loop. Suppose for
contradiction that the system is deadlocked, i.e., no thread ever enters the
CS again. Then after finitely many steps, no thread is making progress in its
doorway either, so every interested thread has reached the spin loop and is
stuck there with a well-defined (label, id) pair.

Among these stuck threads, pick the one with the smallest (label, id) pair;
call it $P_i$. For every other interested thread $k$, either
$\texttt{flag}[k] = \texttt{false}$ (and $P_i$ skips it) or
$(\texttt{label}[k], k) > (\texttt{label}[i], i)$. In both cases, the blocking
conjunction $\texttt{flag}[k] \wedge (\texttt{label}[i], i) >
(\texttt{label}[k], k)$ is false. So $P_i$'s next pass through the spin
loop sees no blocker and enters the CS, contradicting our assumption. In
short: at any point, the smallest-ticketed thread always makes progress.

\section{First-Come-First-Served (FCFS)}

The Bakery Algorithm satisfies FCFS, corresponding to bounded waiting with
$r=0$ in the terminology from Lecture~12. Suppose thread $A$ completes its
doorway before thread $B$ starts its doorway ($D_A \to D_B$). By
Lemma~\ref{lem:label-order}, $(L_A, A) < (L_B, B)$.

We must show $A$ enters the CS no later than $B$ for these acquisitions. Two
cases:

\begin{itemize}[nosep,leftmargin=2.2em]
  \item $A$ has already entered (and possibly even left) the CS by the time $B$
        starts spinning. Then trivially $A$ entered before $B$, and FCFS holds.
  \item $A$ is still in the doorway-or-CS region when $B$ starts spinning.
        Then $\texttt{flag}[A] = \texttt{true}$ continuously throughout
        $B$'s relevant test, and $\texttt{label}[A] = L_A$, so $B$'s check for
        $k=A$ reads $\texttt{flag}[A] = \texttt{true}$ together with
        $(L_B, B) > (L_A, A)$, and $B$ blocks. $A$, on the other hand, either
        already passed all its checks (so it is in the CS) or, when it
        eventually checks $k=B$, sees $(L_A, A) < (L_B, B)$ and is therefore
        not blocked by $B$. So $A$ enters the CS first.
\end{itemize}

In both cases $A$ enters before $B$: no later-arriving thread can overtake an
earlier one even once after the earlier thread has finished its doorway, which
is exactly bounded waiting with $r=0$.

\section{Summary and Looking Ahead}

The Bakery Algorithm achieves $n$-thread mutual exclusion using only shared reads
and writes---no hardware atomics like CAS or TAS are needed.

\begin{ppnote}
\par
Practical limitations to be aware of: (1)~Labels grow without bound and can
overflow; bounded variants such as the Black-White Bakery algorithm exist but
are more complex. (2)~Each lock/unlock requires $O(n)$ reads, making it
impractical for large $n$. (3)~The algorithm uses $O(n)$ shared registers with
a ``single-writer'' property: each thread writes only to its own
\ppinline{flag[i]} and \ppinline{label[i]}, while others only read them.
(4)~The algorithm assumes sequential consistency; weaker models such as TSO
require an explicit store-load fence between each thread's doorway writes
and its spin-loop reads (see the TSO exercise below).
\end{ppnote}

\section{Exercises}

\begin{ppexercises}

  \ppexercise{Why can't we replace \ppinline{flag[i]} with the check
  \ppinline{label[i] != 0}? Construct a concrete counterexample showing how
  mutual exclusion could be violated.

  \medskip
  \noindent\textit{Answer.}
  In the Herlihy-Shavit version (without \ppinline{flag}), the interest
  indicator becomes ``\ppinline{label[k] != 0}.'' The problem is that during
  thread $i$'s doorway, after $i$ has read others' labels but before $i$
  writes its own, $\texttt{label}[i]$ is still $0$, so other threads regard
  $i$ as ``not interested.''

  Concretely, with $i < j$ and all values initialized to $0$:
  \begin{enumerate}[nosep,leftmargin=2.2em]
    \item $P_i$ runs its max-scan to completion, reading
          $\texttt{label}[j] = 0$ (the initial value, since $P_j$ has not
          written yet). $P_i$ has decided $\max = 0$ but defers the
          $\texttt{label}[i] = 1$ write momentarily.
    \item $P_j$ runs its max-scan, reading
          $\texttt{label}[i] = 0$ (still the initial value, since $P_i$ has
          not yet written). $P_j$ writes $\texttt{label}[j] = 1$.
    \item $P_j$ enters the while loop. The check for $k=i$ reads
          $\texttt{label}[i] = 0$, so the blocking conjunction is false; $P_j$
          enters the CS.
    \item $P_i$ now writes $\texttt{label}[i] = 1$ (using the $\max=0$ it
          decided in step~1).
    \item $P_i$ enters the while loop. The check for $k=j$ reads
          $\texttt{label}[j] = 1$ and evaluates $(1, i) > (1, j)$, which is
          \textbf{false} since $i < j$. $P_i$ also enters the CS.
  \end{enumerate}
  Both threads are now in the CS simultaneously, violating ME. The
  \ppinline{flag} array (set \emph{before} the label-write in line~D1) closes
  this gap: any thread seen mid-doorway is already visible to scanners as
  ``interested,'' even before its label is published.}

  \ppexercise{Two threads $T_3$ and $T_7$ both compute label value 5. Which
  enters the CS first, and why?

  \medskip
  \noindent\textit{Answer.}
  The algorithm uses lexicographic comparison on (label, thread\_id). Since both
  have label 5, we compare ids: $3 < 7$, so $(5, 3) < (5, 7)$. Thread $T_3$ has
  the smaller pair and enters first. Thread $T_7$ waits until $T_3$'s flag
  becomes false.}

  \ppexercise{Peterson's algorithm uses a shared \ppinline{victim} variable
  written by all threads. The Bakery algorithm does not. Why does this
  ``single-writer'' property matter, and how strong a guarantee does it buy?

  \medskip
  \noindent\textit{Answer.}
  Two formulations need to be kept distinct.

  In the Herlihy--Shavit formulation used in this lecture, each thread writes
  only to its own \ppinline{flag[i]} and \ppinline{label[i]}; other threads
  only read them. The proof above assumes \emph{atomic} registers, in which
  every read or write appears to take effect at a single instant. Under this
  assumption, the single-writer property is mostly an architectural
  convenience: it sidesteps multi-writer race semantics entirely.

  Lamport's 1974 algorithm (which uses \ppinline{choosing[i]} and resets
  \ppinline{number[i] = 0} on unlock) is correct under a strictly weaker
  hardware model. Lamport proved it works with merely \emph{safe} registers: a
  read that is concurrent with a write to the same location may return any
  value the register can hold, while a read that does not overlap any write
  returns the value most recently written. Single-writer is what makes this
  possible---there is no second writer for a concurrent read to disagree
  with. Peterson's \ppinline{victim} register is multi-writer and so does not
  enjoy this property.

  Bottom line: in the simplified textbook version we use, the single-writer
  property is mainly hygienic; in Lamport's original version it is what
  enables the strongest known register-level correctness result for $n$-thread
  software mutual exclusion.}

  \ppexercise{The proof structures the contradiction around the doorway
  ordering $D_j \to D_i$. Could the two doorways have completed
  ``simultaneously'' under SC?

  \medskip
  \noindent\textit{Answer.}
  No. SC totally orders every memory access in a single global order, so the
  flag-writes and label-writes of $P_i$ and $P_j$ all sit at distinct points
  in that order---one of $D_i$ ends strictly before or strictly after the
  other ends, and likewise for the starts. The proof only needs the strict
  ordering ``$D_j$ ends before $D_i$ begins'' (used by
  Lemma~\ref{lem:label-order}), which is well-defined precisely because SC
  totalizes events. Under weaker models (e.g., release/acquire without
  additional fences), there is no global total order on independent writes,
  and the lemma's appeal to ``the most recent write'' loses its meaning.}

  \ppexercise{Labels in the Bakery algorithm grow without bound. Can we fix
  this by resetting all labels to 0 when the CS is empty?

  \medskip
  \noindent\textit{Answer.}
  This is tempting but dangerous. The problem is that checking ``the CS is
  empty'' and resetting labels is not atomic. While one thread is resetting
  labels, another might be in its doorway reading them, leading to
  inconsistencies. Any fix along these lines requires its own synchronization,
  which defeats the purpose. The standard approach to bounding labels is the
  Black-White Bakery algorithm (Taubenfeld, 2004), which uses one extra shared
  bit to ``color'' rounds and allows labels to wrap around a bounded range. It
  preserves all correctness properties but is more complex.}

  \ppexercise{Does the Bakery algorithm work under the TSO (Total Store Order)
  memory model used by x86 processors? Where exactly do fences need to go?

  \medskip
  \noindent\textit{Answer.}
  Not directly, and the broken ordering is more specific than ``stores aren't
  visible.'' Under TSO, stores by a single thread are made globally visible in
  program order: store-store ordering is preserved for free. So the proof's
  reliance on $\texttt{flag}[i] = \texttt{true}$ becoming visible \emph{before}
  $\texttt{label}[i] = L_i$ does \emph{not} require a fence between (D1) and
  (D2).

  What TSO does \emph{not} preserve is store-load ordering: a thread's own
  pending stores can sit in its store buffer while the same thread issues
  later loads, and other CPUs may keep observing the old value until the
  buffer drains. The dangerous interaction in Bakery is therefore between a
  thread's doorway \emph{writes} and the spin-loop's \emph{reads of other
  threads' flags and labels}. If $P_i$'s store of
  $\texttt{flag}[i] = \texttt{true}$ is still in its buffer when $P_j$
  reaches its check for $k=i$, $P_j$ may read $\texttt{flag}[i] =
  \texttt{false}$, and Step~3 of the ME proof breaks.

  The minimal fix is a single store-load fence (e.g., \ppinline{mfence} on
  x86) after the doorway writes and before the spin loop, forcing each
  thread's doorway writes to drain to memory before it begins reading other
  threads' state. No fence is needed between (D1) and (D2).}

  \ppexercise{Every thread scans all $n$ labels in both the doorway (to compute
  the max) and the while loop. Is this $O(n)$ cost per operation inherent? What
  is the relevant lower bound?

  \medskip
  \noindent\textit{Answer.}
  Two related lower bounds need to be kept apart.
  \begin{itemize}[nosep,leftmargin=2.2em]
    \item \textbf{Space.} Burns and Lynch (1980) showed that any $n$-thread
          mutual exclusion algorithm using only reads and writes to shared
          registers must use at least $n$ distinct shared locations. So
          Bakery's $O(n)$ shared registers are space-optimal up to a
          constant; \emph{this is a statement about how many registers
          exist, not how many a single acquisition reads}.
    \item \textbf{Per-acquisition reads.} Bakery's $O(n)$ reads per
          \ppinline{lock()}\,/\,\ppinline{unlock()} pair are structural to
          the label-comparison strategy: a thread cannot determine its
          relative position without inspecting every other thread's label
          and flag. Subsequent work (e.g., Lamport's $1$-bit and fast mutual
          exclusion algorithms; Yang and Anderson's local-spinning locks)
          shows that the cost can be reduced in restricted regimes (e.g.,
          contention-free or local-spin) but not eliminated in the worst
          case under reads/writes alone.
  \end{itemize}
  The practical takeaway is that read-only solutions are unavoidably linear
  per operation in the worst case, which is why production systems lean on
  hardware primitives such as TAS/CAS (Lecture~11) to bring the
  contention-free cost down to $O(1)$.}

\end{ppexercises}

\section*{Contributions}

This scribe was prepared by \textbf{Group 17} as part of the class scribe
rotation for \textit{Introduction to Parallel \& Distributed Programming}.

\bigskip
\renewcommand{\arraystretch}{1.4}
\begin{tabularx}{\textwidth}{|p{0.30\textwidth}|p{0.20\textwidth}|X|}
\hline
\textbf{Name} & \textbf{Entry Number} & \textbf{Contribution} \\
\hline
Arnav Raj & 2022CS51652 & Overall structure, proof verification against
literature, exercises and answers, final editing \\
\hline
Kushagar Garg & 2022CS11088 & Exercise and discussion questions, sketching the
overall layout, editing in different sections \\
\hline
Irfan Hussain & 2022CS51146 & Reviewed the final transcript and made
corrections to improve clarity and readability \\
\hline
Meda Sadvik & 2023CS11097 & Assisted with minor formatting adjustments and
reviewed selected sections for consistency \\
\hline
Eeshan Anshuman Yadav & 2023CS50917 & Contributed to initial reading of lecture
material and provided brief feedback on explanations \\
\hline
\end{tabularx}

\bigskip
\noindent All group members reviewed the final draft for correctness and
clarity. The mutual exclusion proof was cross-checked against Lamport's
original 1974 paper, the Herlihy \& Shavit textbook (\textit{The Art of
Multiprocessor Programming}), Garg's textbook (\textit{Concurrent and
Distributed Computing in Java}), and Shankar's lecture notes (UMD CMSC~412).